% Parsimony -- Project Report (short version)
% Compile on Overleaf: paste this file, click Recompile. Engine: pdfLaTeX.
\documentclass[11pt,a4paper]{article}

\usepackage[margin=2.2cm]{geometry}
\usepackage{times}
\usepackage[T1]{fontenc}
\usepackage[utf8]{inputenc}
\usepackage{booktabs}
\usepackage{array}
\usepackage{tabularx}
\usepackage{xcolor}
\usepackage[most]{tcolorbox}
\usepackage{tikz}
\usetikzlibrary{positioning,arrows.meta}
\usepackage{enumitem}
\usepackage{fancyhdr}
\usepackage[hidelinks]{hyperref}

\definecolor{ink}{HTML}{16181D}
\definecolor{accent}{HTML}{7A2518}
\definecolor{soft}{HTML}{FBFAF7}

\pagestyle{fancy}
\fancyhf{}
\lhead{\small Parsimony: Token-Efficient LLM Interaction}
\rhead{\small VIT University}
\cfoot{\small\thepage}
\renewcommand{\headrulewidth}{0.4pt}

\setlist{nosep,leftmargin=1.3em}
\newcolumntype{L}[1]{>{\raggedright\arraybackslash}p{#1}}

% A titled box for each research gap
\newtcolorbox{gap}[1]{
  colback=soft, colframe=ink, coltitle=white,
  fonttitle=\bfseries\small, title={#1},
  boxrule=0.7pt, arc=2pt, left=6pt, right=6pt, top=5pt, bottom=5pt,
  breakable
}
\newcommand{\lbl}[1]{\textcolor{accent}{\textbf{#1}}}

\begin{document}

% ------------------------------------------------------------------ cover
\thispagestyle{empty}
\begin{center}
  \vspace*{1.2cm}
  {\large\bfseries VIT UNIVERSITY}\\[2pt]
  {\small Vellore Institute of Technology}\\[1pt]
  {\small Vellore, Tamil Nadu, India}

  \vspace{2.6cm}
  {\huge\bfseries Token-Efficient LLM Interaction\\[4pt] on CPU-Only Hardware}\\[10pt]
  {\large\itshape Parsimony --- A Stacked, Self-Improving\\ Optimisation Layer for Small Language Models}

  \vspace{1.6cm}
  {\normalsize Natural Language Processing Project}

  \vspace{1.4cm}
  {\small\textcolor{accent}{\textbf{TEAM MEMBERS}}}\\[8pt]
  \begin{tabular}{ll}
    Arrsh Tripathi & 23BCI0191 \\
    Alok Singh     & 23BCI0158 \\
    Dev Singh      & 23BCE0794 \\
  \end{tabular}

  \vspace{1.8cm}
  {\small B.Tech \quad$\cdot$\quad Course Code: BCSE497J \quad$\cdot$\quad Project I}\\[4pt]
  {\small Faculty Guide: Dr Sathya K}\\[4pt]
  {\small September 2026}
\end{center}
\newpage

% --------------------------------------------------------------- abstract
\section*{Abstract}
Language models charge by the token. On a laptop with no graphics card, they also
spend nearly all of their time \emph{reading} the question rather than
\emph{writing} the answer. Many methods exist to cut that cost --- shortening the
prompt, reusing old answers, trimming chat history, capping the reply length ---
but each one is published on its own. Nobody has run them together and measured
what happens.

We read \textbf{40 papers}, list what each one cannot do in our setting, and pull
out \textbf{six research gaps}. We then build \textbf{Parsimony}: a layer that sits
between the app and a small local model. It has seven optimisation modules and one
always-on safety check. Every module only \emph{suggests} a change; one controller
decides whether to accept it. Because each module has an on/off switch, the whole
system doubles as an experiment.

Measured on 151 conversations and 263 requests against \texttt{qwen2.5:1.5b-instruct}
running on CPU only, the full stack removes \textbf{33.9\%} of tokens with
\textbf{no drop in answer quality} (92.5\%~$\rightarrow$~97.5\% on 40 test questions,
zero answers made worse). Three results are new: savings \textbf{do not add up} when
modules are stacked; reading the prompt is \textbf{92--99\%} of the time on CPU
($\approx$8.5\,ms per input token); and the cache similarity thresholds published as
safe are \textbf{unsafe} here, so a cheap verifier --- not a threshold --- is what
makes answer reuse safe (wrong-answer rate 26.7\%~$\rightarrow$~0.0\%).

% ------------------------------------------------------------- lit review
\section{Literature Review}

Forty papers were read across six strands. For each we note the method, the
limitation \emph{for our setting}, and what it measures. These are not criticisms of
the papers on their own terms --- most target datacentre servers or a single
technique. Our setting is simply the one they do not target.

\subsection{Limitations of the 40 papers}

\begin{tabularx}{\textwidth}{@{}l L{3.6cm} X l@{}}
\toprule
\textbf{Ref} & \textbf{Method} & \textbf{Limitation for us} & \textbf{Metrics} \\
\midrule
\multicolumn{4}{@{}l}{\textit{Strand A --- Prompt compression}}\\
[1] & Selective Context & Needs a second model; never tested with a cache & latency, score \\
[2] & LLMLingua & 20$\times$ compression, but GPU and alone & ratio, score \\
[3] & LongLLMLingua & Very long contexts; a laptop rarely gets there & speedup, score \\
[4] & LLMLingua-2 & Faithfulness claimed, never attacked & ratio, faithfulness \\
[5] & Evaluator heads & Needs the model's internals & latency, accuracy \\
[6] & SCOPE & Adds generation cost to the same request & ratio, quality \\
[7] & Gisting & Needs fine-tuning & ratio, score \\
[8] & UltraGist & Needs training & ratio, accuracy \\
[9] & ATACompressor & Assumes the task label is known & ratio, score \\
[10] & SARA & Needs a retrieval corpus & retrieval quality \\
[11] & PCToolkit & A toolkit, not a method; no cross-family mixing & comparisons \\
[12] & Empirical study & \textbf{Names} two failure modes; no run-time guard & hallucination rate \\
[13] & Information preservation & Diagnoses a 30--50 point drop; no fix & groundedness \\
[14] & Financial fidelity & Shows fluent but decision-changing output & agreement \\
[15] & Compression paradox & \textbf{Does not always save energy}; hosted APIs only & energy \\
\addlinespace
\multicolumn{4}{@{}l}{\textit{Strand B --- Semantic caching}}\\
[16] & GPTCache & One threshold, no verification, never attacked & hit rate, cost \\
[17] & GPT Semantic Cache & Same design, different storage & hit rate, latency \\
[18] & MeanCache & Hit rate \emph{inside} a chat not studied & hit rate, privacy \\
[19] & ContextCache & Closest to us; no false-hit rate reported & hit rate \\
[20] & Generative caching & Rewrites instead of reusing; costs more & hit rate, quality \\
[21] & Key-collision attack & Treats it as a deliberate attack; we face ordinary pairs & attack rate \\
[22] & SAFE-CACHE & Right conclusion, but a heavy redesign & robustness \\
\addlinespace
\multicolumn{4}{@{}l}{\textit{Strand C --- KV cache and prefix reuse}}\\
[23] & PagedAttention / vLLM & GPU serving; app-side cost not priced & throughput \\
[24] & Prefix caching & States the rule; never prices breaking it & reuse rate \\
[25] & ChunkAttention & Kernel level; many users assumed & throughput \\
[26] & Sparse prefix caching & Tied to one architecture & throughput \\
[27] & Multi-segment attention & Server side; no app guidance & latency, memory \\
\addlinespace
\multicolumn{4}{@{}l}{\textit{Strand D --- Where text sits in the prompt}}\\
[28] & Lost in the Middle & Says reorder by relevance --- which \textbf{breaks prefix reuse} & accuracy by position \\
[29] & Adaptive Focus Memory & Rules of thumb; reorder cost unmeasured & token usage \\
[30] & Multi-turn survey & Survey; nothing composed & mixed \\
[31] & Memory survey & Taxonomy, not measurement & mixed \\
\addlinespace
\multicolumn{4}{@{}l}{\textit{Strand E --- Routing between models}}\\
[32] & FrugalGPT & 98\% saved \emph{against paid APIs}; assumes a big model exists & cost, accuracy \\
[33] & RouteLLM & Assumes the bigger model is better --- false at 1.5B vs 3B here & cost, scores \\
[34] & Semantic router & Datacentre assumption & latency, cost \\
[35] & UCCI & Needs calibrated uncertainty; costly on CPU & cost, accuracy \\
[36] & On-device routing & Routes device$\leftrightarrow$cloud, not inside one stack & accuracy \\
[37] & Cluster, route, escalate & No quality guarantee on the cheap path & cost, quality \\
\addlinespace
\multicolumn{4}{@{}l}{\textit{Strand F --- Controlling answer length}}\\
[38] & Length control encoding & \textbf{Needs fine-tuning} & length adherence \\
[39] & BudgetThinker & Needs training with control tokens & budget, accuracy \\
[40] & Strict length limits & Measures the damage; no budgeting scheme & accuracy vs length \\
\bottomrule
\end{tabularx}

\subsection{What keeps repeating}

\begin{tabularx}{\textwidth}{@{}l X l@{}}
\toprule
\textbf{\#} & \textbf{Limitation that recurs} & \textbf{Papers} \\
\midrule
L1 & Tested alone, never combined with another method & [1]--[11], [16]--[20], [38]--[40] \\
L2 & Assumes a GPU server, where costs are the opposite of a laptop & [2], [3], [23]--[27], [32]--[35] \\
L3 & Failures are described but not \emph{prevented} while running & [12]--[14], [21], [22] \\
L4 & Settings quoted as universal, tuned only on GPT-class models & [2], [4], [16], [17], [19] \\
L5 & Needs training, fine-tuning, or the model's internals & [5], [7], [8], [38], [39] \\
\bottomrule
\end{tabularx}

\subsection{Papers closest to our work}

\textbf{[28] Lost in the Middle} and \textbf{[24] prefix caching} contradict each
other. The first says move important text to the front. The second says the front
must never change. \emph{Doing the first destroys the second}, and neither paper
mentions the other. We price this at about 80$\times$ (\S6.2).

\textbf{[16] GPTCache} set the pattern every cache follows: compare, then serve if
similar enough. Our attack set shows the threshold is the wrong place to put the
safety (\S6.3). \textbf{[22] SAFE-CACHE} reaches the same conclusion but answers it
with a redesign; we answer it with four cheap checks.

\textbf{[12] Empirical Study on Prompt Compression} names the two ways compression
breaks an answer and counts how often it happens --- but does not stop it. Our
module M8 is that stop.

% -------------------------------------------------------------------- gaps
\section{Research Gaps Identified}

\subsection{Gap 1: The techniques have never been stacked}
\begin{gap}{Gap 1 --- Isolated Modules, Unknown Interactions}
\lbl{Current Problem:} Compression, caching, history trimming and routing each have
their own literature, benchmarks and tuned settings. No published study runs all of
them in one live pipeline, so nobody knows whether their savings add up.

\lbl{Research Gap:} Do the savings add up, or do the modules eat each other's lunch?
A compressor that rewrites a query changes what the cache sees. A trimmer that drops
turns changes what the compressor has left to remove. Nobody has measured the sign or
size of these interactions.

\lbl{Our Solution:} A complete $2^4$ factorial ablation over the four headline
modules, reported with confidence intervals and an analysis of variance on the
interaction terms.
\end{gap}
\textbf{Answered:} savings fall short by \textbf{1.63\,pp}, 95\% CI $[-0.02, +3.23]$,
and both real interaction terms are negative (\S6.1).

\subsection{Gap 2: The output side of the bill is ignored}
\begin{gap}{Gap 2 --- Everyone Optimises the Prompt, Nobody the Answer}
\lbl{Current Problem:} Compression, caching and pruning all shrink the input. That
made sense for paid APIs, where the prompt is what is billed. On a consumer CPU the
economics change: reading the prompt is a parallel operation, writing the answer is
strictly one token at a time.

\lbl{Recent Evidence:} A length hint can cut output tokens by 74--86\% on chatty
models, but the same hint hurts models whose answers are already short. Length must
be set per question.

\lbl{Research Gap:} No study reports the read/write cost split for CPU-hosted small
models, and no optimisation stack contains an output-side module at all.

\lbl{Our Solution:} Measure time-to-first-token and time-per-output-token separately,
and add an adaptive output budgeter with a streaming early stop as its own module.
\end{gap}
\textbf{Answered, and the expected direction was wrong:} reading the prompt is
\textbf{92--99\%} of the time (\S6.2). The budgeter is still the biggest single token
saver (+13.44\,pp), because it is the only module touching the output at all.

\subsection{Gap 3: Compression and caching have never met}
\begin{gap}{Gap 3 --- The Compression--Cache Interaction Is Unmeasured}
\lbl{Current Problem:} Semantic caches produce false hits: two questions that look
alike but mean different things get matched, and the user is given the answer to a
question they did not ask.

\lbl{Recent Discovery:} Nobody has asked what compression does to that. It cuts both
ways. Compression strips small words, and small words are often exactly what tells
two similar questions apart --- so false hits could get worse. But compression also
makes wording uniform, so two ways of asking the same thing match --- so true hits
could get better.

\lbl{Research Gap:} Does compression raise or lower the false-hit rate, and does the
safe threshold move?

\lbl{Our Solution:} A purpose-built set of near-identical question pairs differing in
exactly one important word, swept across thresholds with compression on and off.
\end{gap}
\textbf{Answered, and the question dissolves:} no threshold separates the two classes
at all, so the answer is a verifier, not a different threshold (\S6.3).

\subsection{Gap 4: Fewer tokens can cost more time}
\begin{gap}{Gap 4 --- Optimisation That Defeats the Key-Value Cache}
\lbl{Current Problem:} Local runtimes reuse their key-value cache for any prompt
opening identical to the last request, so a stable opening is processed once and then
skipped. This is the biggest speed win on CPU, and a token counter cannot see it.

\lbl{Recent Discovery:} Every technique in the literature rewrites the prompt. A
compressor that changes the first sentence, or a trimmer that drops an early turn,
throws that reuse away. The system then sends fewer tokens and takes longer.

\lbl{Research Gap:} No published work measures token savings and prefix reuse
together. Token count is treated as a stand-in for cost, and on local hardware that
stand-in can invert.

\lbl{Our Solution:} A prompt assembler that pins unchanging content to the front and
keeps all rewriting at the end, plus a direct measurement of where more compression
starts costing time.
\end{gap}
\textbf{Answered:} the same content, 0.5\% apart in tokens, is \textbf{80$\times$
apart in cost} --- 212\,ms against 18{,}914\,ms (\S6.2).

\subsection{Gap 5: Thresholds tuned on models that do not break}
\begin{gap}{Gap 5 --- Small Models Are More Brittle Than the Tuning Set}
\lbl{Current Problem:} The compression ratios and cache thresholds quoted as safe
were all calibrated against GPT-3.5 and GPT-4 class models.

\lbl{Recent Discovery:} Small open models change their answer across rewordings that
mean the same thing, while GPT-class models rarely do. A compressed prompt is, to the
model, an aggressive rewording.

\lbl{Research Gap:} Do published safe settings transfer downward, or does the safe
ratio for a 1B model sit far below the published figure?

\lbl{Our Solution:} Re-run the standard protocol on local models and publish a
calibration table per model instead of one universal number.
\end{gap}
\textbf{Answered:} the published 0.85--0.92 cache range is \textbf{unsafe here} --- it
accepts a negation pair at 0.924 and returns the opposite answer (\S6.3). Reduction
\emph{ratios}, though, do transfer within 0.1\,pp (\S6.5).

\subsection{Gap 6: Every cache in the literature starts empty}
\begin{gap}{Gap 6 --- The Unused Signal in the User's Own Conversations}
\lbl{Current Problem:} Every caching result is reported on a warm cache. On day one
the cache is empty, the hit rate is zero, and the user pays full price for exactly the
questions they have asked many times before in some other chat window.

\lbl{Recent Discovery:} MeanCache showed that about a third of real queries closely
resemble an earlier one, and that keeping the cache on the user's own device is both
workable and more private --- but it still starts from scratch.

\lbl{Research Gap:} A user's own chat history is a large, free, perfectly on-topic
corpus that no optimisation system uses. Nothing in the literature mines it to warm up
a cache or to learn the phrases that user always writes and never needs.

\lbl{Our Solution:} An offline policy learner that replays exported chat logs,
re-scores each turn under candidate policies, and produces a pre-filled cache, a
personal redundancy list and a context digest --- all locally, with nothing leaving the
machine.
\end{gap}
\textbf{Answered, and made conditional:} mined policy is worth \textbf{+0.00\,pp when
nothing repeats and +17.83\,pp at 57\% repetition} (\S6.5). ``Self-improving'' is a
property of the traffic, not of the module.

% ------------------------------------------------------------- RQ and PS
\section{Research Questions and Problem Statement}

\textbf{RQ1.} When several token-saving techniques run in one pipeline, do their
savings add up --- and if not, by how much do they fall short? \emph{(Gap 1)}

\textbf{RQ2.} On CPU-only hardware, which half of a request costs more, and what is
one input token worth in milliseconds? \emph{(Gaps 2, 4)}

\textbf{RQ3.} Can answer reuse be made safe against near-identical, opposite-meaning
questions without running a second model? \emph{(Gaps 3, 4)}

\textbf{RQ4.} Do settings tuned in one configuration transfer to another, or to
conversations never seen before? \emph{(Gaps 5, 6)}

\subsection*{Problem statement}
\begin{tcolorbox}[colback=soft,colframe=accent,boxrule=0pt,leftrule=2.5pt,arc=0pt]
Build a middleware layer that measurably cuts the token cost of using a small language
model on CPU-only consumer hardware, \textbf{without changing the meaning of any
request}, and instrument it so that the contribution of every module, and every
interaction between modules, can be measured separately and reproduced independently.
\end{tcolorbox}
The emphasis is deliberate. A system that saves tokens by quietly damaging requests is
not a contribution; safety is the constraint every optimisation must work under.

% ------------------------------------------------------- contributions
\section{Contributions}

\begin{enumerate}[label=\textbf{C\arabic*.}]
\item \textbf{The shortfall, measured.} A full $2^4$ factorial over compressor,
cache, history manager and output budgeter. Savings do not add up: the shortfall is
\textbf{1.63\,pp}, 95\% CI $[-0.02,+3.23]$. Its size depends on the configuration ---
a better encoder moved it from 2.53 to 1.63\,pp. \emph{(RQ1, Gap 1)}

\item \textbf{The CPU cost structure, and the price of prompt order.} Reading the
prompt is \textbf{92--99\%} of the time, straight-line at \textbf{$\approx$8.5\,ms per
input token}. This turns every token result into seconds. It also prices a conflict
the literature never names: one changing word at the very start of a prompt costs
\textbf{212\,ms $\rightarrow$ 18{,}914\,ms} for the same content. \emph{(RQ2, Gaps 2, 4)}

\item \textbf{A verifier that makes reuse safe without a second model.} 45 attack
pairs one word apart, plus 45 controls. The negation pair scores \textbf{0.924 ---
above every genuine reworded pair}, so no threshold can separate them. Four cheap set
comparisons take the wrong-answer rate from \textbf{26.7\% to 0.0\%}. \emph{(RQ3,
Gaps 3, 4)}

\item \textbf{Transfer, tested both ways.} Against a second vocabulary the reduction
percentages transfer within 0.1\,pp with the same module ranking, but the
\emph{explanations} do not --- one of our own two explanations turned out to be true of
only one tokenizer. Mined policy transfers as a function of repetition:
\textbf{+0.00\,pp at 0\%, +17.83\,pp at 57\%}, with zero safety violations at any
level. \emph{(RQ4, Gaps 5, 6)}
\end{enumerate}

% ---------------------------------------------------------------- method
\section{Proposed Method}

\subsection{The one design decision everything follows from}
Every module \textbf{suggests} a change; one controller \textbf{applies} it. A
suggestion is a description of an edit, not the edit itself. Three things follow that
no module has to implement:

\begin{enumerate}
\item \textbf{Clean on/off switches.} Turning a module off really removes its effect,
so an experiment cell means what it claims.
\item \textbf{One safety checkpoint.} Every edit passes one gate, so the safety check
is not written seven times.
\item \textbf{Undo is free.} A refused suggestion is simply never applied.
\end{enumerate}

Stage order is \textbf{a setting, not code}. Gap 3 is only answerable because the
cache can be moved before or after the compressor without editing the pipeline.

\subsection{The pipeline, tier by tier}

\begin{tabularx}{\textwidth}{@{}c L{3.3cm} X l@{}}
\toprule
\textbf{\#} & \textbf{Tier} & \textbf{What it does} & \textbf{Effect} \\
\midrule
1 & M6a Deterministic router & Works out maths and date questions exactly, without the model & 23 tokens $\rightarrow$ \textbf{0} \\
2 & M2 Semantic cache & Looks for an earlier answer; borderline matches are checked by four comparisons before reuse & prompt $\rightarrow$ stored answer \\
3 & M3a History selector & Keeps only the earlier turns that matter to this question & 8 turns $\rightarrow$ 6 \\
4 & M3b History arranger & Decides the order of the kept turns & order changes cost \\
5 & M1a Tier 1 & Removes greetings, politeness and formatting & 28 $\rightarrow$ 17 tokens \\
6 & M1b Tier 2 & Deletes a sentence that repeats a fact already said & 36 $\rightarrow$ 26 tokens \\
7 & M1c Tier 3 & Shortens wordy phrasing, and rejects any edit that does not actually save tokens & rejects no-gain edits \\
8 & M4 Prefix-stable assembler & Puts unchanging text first so the model can reuse its cache & 212\,ms vs 18{,}914\,ms \\
9 & M5 Output budgeter & Sets the answer length by question type, plus an early stop & 48 \ldots 640 tokens \\
10 & M6b Escalation router & Decides if a bigger model is needed & calibrated, off by default \\
\midrule
--- & M8 Fidelity gate & \textbf{Checks every suggestion} from every stage; refuses anything that would drop a number, name, ``not'' or qualifier & 26.7\% $\rightarrow$ 0.0\% \\
--- & M7 Policy learner & Offline: mines past chats into a bundle that pre-fills the cache & +17.83\,pp at 57\% repetition \\
\bottomrule
\end{tabularx}

\subsection{Experimental protocol}
\begin{tabularx}{\textwidth}{@{}l X@{}}
\toprule
Design & $2^4$ full factorial over M1 $\times$ M2 $\times$ M3 $\times$ M5, 17 cells \\
Corpus & 151 conversations, 263 requests; 45 attack pairs + 45 controls; 40 test questions \\
Model & \texttt{qwen2.5:1.5b-instruct}, 4-bit, via Ollama, CPU only, offline \\
Hardware & AMD Ryzen 7 5800HS, 16\,GB RAM, no GPU used \\
Statistics & Bootstrap 95\% intervals, partial $\eta^2$, two-way ANOVA, exact McNemar \\
Reproduction & \texttt{python reproduce.py} --- every table regenerates in $\approx$100\,s \\
Verification & 719 automated tests, including one that fails the build if a layer imports upward \\
\bottomrule
\end{tabularx}

\newpage
% ---------------------------------------------------------- architecture
\section{System Architecture}

This section lists what is actually inside the system: its layers, its parts, and the
record it leaves behind.

\subsection{The request path}
\begin{center}
\begin{tikzpicture}[
  font=\scriptsize,
  box/.style={draw, rounded corners=2pt, minimum width=7.4cm, minimum height=0.52cm,
              align=left, inner xsep=6pt, text width=7.1cm},
  cap/.style={draw, fill=black!6, rounded corners=2pt, minimum width=7.4cm,
              minimum height=0.5cm, align=center, font=\scriptsize\bfseries},
  ar/.style={-latex, gray}
]
\node[cap] (in) {USER QUESTION};
\node[box, below=2.6pt of in, draw=green!40!black] (m6a)
  {\textbf{M6a} Deterministic router \hfill 23 tokens $\rightarrow$ 0};
\node[box, below=2.6pt of m6a, draw=green!40!black] (m2)
  {\textbf{M2} Semantic cache \hfill reuse a stored answer};
\node[box, below=2.6pt of m2, draw=blue!50!black] (m3a)
  {\textbf{M3a} History selector \hfill 8 turns $\rightarrow$ 6};
\node[box, below=2.6pt of m3a, draw=blue!50!black] (m3b)
  {\textbf{M3b} History arranger \hfill order changes cost};
\node[box, below=2.6pt of m3b, draw=violet!60!black] (m1a)
  {\textbf{M1a} Tier 1 --- boilerplate \hfill 28 $\rightarrow$ 17 tokens};
\node[box, below=2.6pt of m1a, draw=violet!60!black] (m1b)
  {\textbf{M1b} Tier 2 --- repeated facts \hfill 36 $\rightarrow$ 26 tokens};
\node[box, below=2.6pt of m1b, draw=violet!60!black] (m1c)
  {\textbf{M1c} Tier 3 --- word level \hfill rejects no-gain edits};
\node[box, below=2.6pt of m1c, draw=red!50!black] (m4)
  {\textbf{M4} Prefix-stable assembler \hfill 212\,ms vs 18{,}914\,ms};
\node[box, below=2.6pt of m4, draw=orange!70!black] (m5)
  {\textbf{M5} Output budgeter \hfill 48 \ldots 640 tokens};
\node[box, below=2.6pt of m5, draw=black!60] (m6b)
  {\textbf{M6b} Escalation router \hfill off by default};
\node[cap, below=2.6pt of m6b] (out) {PROMPT SENT TO THE MODEL};

\foreach \a/\b in {in/m6a, m6a/m2, m2/m3a, m3a/m3b, m3b/m1a, m1a/m1b,
                   m1b/m1c, m1c/m4, m4/m5, m5/m6b, m6b/out}
  \draw[ar] (\a) -- (\b);

\node[draw=yellow!50!black, fill=yellow!8, rounded corners=3pt, text width=2.7cm,
      align=center, right=8pt of m1a, minimum height=5.4cm, font=\scriptsize] (gate)
  {\textbf{M8 $\cdot$ FIDELITY GATE}\\[4pt]
   Checks \emph{every} suggestion above before it is applied.\\[4pt]
   Refuses any edit that would drop a number, name, ``not'' or qualifier.\\[5pt]
   \textbf{Wrong answers}\\ \textbf{26.7\% $\rightarrow$ 0.0\%}};
\draw[-latex, yellow!50!black] (gate.west) -- ++(-0.35,0);
\end{tikzpicture}

{\footnotesize\textit{Figure 1.} The request path. Each tier may answer outright,
suggest an edit, or do nothing --- and records which. Every suggestion passes the gate
before it is applied.}
\end{center}

\subsection{What is in the architecture}

The system is built in six layers. A layer may use only the layers below it, and an
automated test fails the build if that is broken.

\begin{tabularx}{\textwidth}{@{}l l X@{}}
\toprule
\textbf{Layer} & \textbf{Name} & \textbf{What it holds} \\
\midrule
L0 & Contracts & The data types only: \texttt{Request}, \texttt{Turn}, \texttt{Proposal}, \texttt{LedgerRow}. No logic. \\
L1 & Infrastructure & Tokeniser (Qwen2.5, 151{,}665 entries), two encoders, text clean-up, model providers, result memoisation. \\
L2 & Modules & M1--M8. Each reads its inputs and returns a \emph{suggestion}; none changes the request itself. \\
L3 & Orchestrator & Runs the stages in the configured order, sends every suggestion to the gate, applies or refuses, writes one log row per stage. \\
L4 & Evaluation & Factorial sweep runner, quality measures, statistics, threshold calibration, cross-vocabulary study, latency instrument, learning study. \\
L5 & Surfaces & The command line: \texttt{run}, \texttt{sweep}, \texttt{compare}, \texttt{tour}, \texttt{ask}. \\
\bottomrule
\end{tabularx}

\subsection{The eight modules}
\begin{itemize}
\item \textbf{M1 Compressor} --- three tiers, sentence-aware, with a guard that rejects any edit not actually saving tokens.
\item \textbf{M2 Cache} --- exact match plus a similarity tier with three zones (accept / check / reject) and a four-check verifier.
\item \textbf{M3 History manager} --- four ways to choose turns, and a separate stage to order them, so \emph{what} is kept and \emph{where} it goes can be switched on and off apart.
\item \textbf{M4 Assembler} --- builds the prompt unchanging-part-first, and reports the reuse percentage directly.
\item \textbf{M5 Output budgeter} --- an answer length per question type, plus an early stop when the model starts repeating itself.
\item \textbf{M6 Router} --- one tier that answers maths and dates using zero model tokens, and one that decides on a bigger model.
\item \textbf{M7 Policy learner} --- offline replay of past chats into a bundle that pre-fills the cache.
\item \textbf{M8 Fidelity gate} --- always on; checks numbers, names, negations, qualifiers and emptiness on every suggestion.
\end{itemize}

\subsection{Parts that cut across everything}
\begin{tabularx}{\textwidth}{@{}l X@{}}
\toprule
\textbf{Ledger} & One row per stage per request: module, outcome, tokens before and after, time taken, reason, gate event, model fingerprint. \textbf{A stage that did nothing still writes a row} --- an invisible stage cannot be audited. \\
\textbf{Fidelity gate} & The single safety checkpoint, before anything is applied. \\
\textbf{Provider layer} & Ollama and a deterministic stand-in. A run that asks for the real model and cannot reach it \textbf{stops} rather than quietly using the fake one. \\
\textbf{Memoisation} & Identical generations are not paid for twice during a sweep. \\
\textbf{Order validator} & Checks the configured stage order against a dependency graph before running. \\
\textbf{Corpus} & 151 conversations, 45 attack pairs, 45 controls, 40 test questions. \\
\textbf{Tests} & 719, including the layering test. \\
\bottomrule
\end{tabularx}

\subsection{Technology stack}
Python 3.11; Ollama running \texttt{qwen2.5:1.5b-instruct} (4-bit) on CPU;
\texttt{tokenizers} for the model's own vocabulary; \texttt{numpy} for maths;
\texttt{typer} and \texttt{rich} for the demo; \texttt{pytest} for the tests.
Deliberately absent: PyTorch, any GPU requirement, any paid service.

\newpage
% --------------------------------------------------------------- results
\section{Results and Discussion}

Every figure regenerates from raw logs with one command in about 100 seconds.
\textbf{719 automated tests pass.}

\subsection{The headline: savings do not add up}
Full $2^4$ factorial, 151 conversations, 263 requests, 17 cells.

\begin{center}
\begin{tabular}{@{}lrr@{}}
\toprule
\textbf{Effect} & \textbf{Estimate} & \textbf{Partial $\eta^2$} \\
\midrule
M5 output budgeter & +13.44\,pp & 0.556 \\
M3 history manager & +11.82\,pp & 0.430 \\
M2 semantic cache  & +1.94\,pp  & 0.012 \\
M1 compressor      & +0.23\,pp  & 0.000 \\
M3 $\times$ M5 interaction & \textbf{$-$0.70\,pp} & 0.002 \\
\bottomrule
\end{tabular}
\end{center}

Full stack: \textbf{+33.9\%} total token reduction. Both real interaction terms are
negative and both involve M5 --- trimming history and shortening the answer are cutting
the same conversation. \textbf{Shortfall: 1.63\,pp, 95\% CI $[-0.02,+3.23]$.}

And the sharper result: the shortfall depends on the configuration.

\begin{center}
\begin{tabular}{@{}lrrl@{}}
\toprule
\textbf{Encoder} & \textbf{M2 effect} & \textbf{M3$\times$M5} & \textbf{Shortfall} \\
\midrule
\texttt{hashing-v1} & +1.61\,pp & $-$1.14 & 2.53\,pp, $[+0.93,+3.99]$ \\
\texttt{content-v1} (default) & +1.94\,pp & $-$0.70 & 1.63\,pp, $[-0.02,+3.23]$ \\
\bottomrule
\end{tabular}
\end{center}

A better encoder made the cache hit more often, so it overlapped its neighbours less.
The weaker encoder was \textbf{not} put back to protect the interval: choosing a
component you know is worse because it gives a nicer number is exactly the failure this
project is written against.

\subsection{Reading the prompt dominates, and prompt order has a price}
Measured on \texttt{qwen2.5:1.5b-instruct}, Ryzen 7 5800HS, using the server's own
counters rather than a stopwatch.

\begin{center}
\begin{tabular}{@{}rrrr@{}}
\toprule
\textbf{Input tokens} & \textbf{Reading} & \textbf{Writing} & \textbf{Reading share} \\
\midrule
146 & 1{,}360\,ms & 122\,ms & 91.7\% \\
474 & 3{,}902\,ms & 165\,ms & 95.9\% \\
1{,}338 & 11{,}609\,ms & 136\,ms & 98.8\% \\
\bottomrule
\end{tabular}
\end{center}

\textbf{$\approx$8.5\,ms per input token.} The 11{,}836 tokens the full stack saves are
\textbf{100.6 seconds} across the corpus --- 383\,ms per request.

\begin{center}
\begin{tabular}{@{}lrrr@{}}
\toprule
\textbf{Prompt arrangement} & \textbf{Tokens} & \textbf{Steady-state read} & \textbf{Reuse} \\
\midrule
Unchanging front (M4) & 1{,}490 & \textbf{212\,ms} & \textbf{98.5\%} \\
Changing front & 1{,}497 & \textbf{18{,}914\,ms} & 0.5\% \\
\bottomrule
\end{tabular}
\end{center}

Same content, 0.5\% apart in tokens, about \textbf{80$\times$ apart in cost}. Every
metric in the compression literature scores these two the same.

\subsection{The verifier, not the threshold}
45 attack pairs plus 45 controls.

\begin{center}
\begin{tabular}{@{}rl@{}}
\toprule
\textbf{Similarity} & \textbf{Pair} \\
\midrule
\textbf{0.924} & ``Is it safe to mix bleach and vinegar?'' / ``Is it \textbf{not} safe\ldots'' \\
0.869 & ``capital of \textbf{Australia}'' / ``capital of \textbf{Austria}'' \\
0.653 & ``What causes rain?'' / ``What causes rainfall to occur?'' \\
\bottomrule
\end{tabular}
\end{center}

The attack pair scores \emph{above every genuine reworded pair}, so the published
0.85--0.92 range would accept it and return the opposite answer.

\begin{center}
\begin{tabular}{@{}lrr@{}}
\toprule
\textbf{Difference} & \textbf{Before} & \textbf{After the verifier} \\
\midrule
qualifier & 78\% & \textbf{0\%} \\
negation  & 31\% & \textbf{0\%} \\
name      & 11\% & \textbf{0\%} \\
\textbf{overall} & \textbf{26.7\%} & \textbf{0.0\%} \\
\bottomrule
\end{tabular}
\end{center}

\subsection{Quality is not the price}
On 40 test questions with the real model: baseline \textbf{37/40 (92.5\%)}, full stack
\textbf{39/40 (97.5\%)}. Compared item by item, \textbf{no answer that was right before
became wrong}. The two gains are maths questions handled by M6a without the model; two
differing pairs is not statistically significant (McNemar $p = 0.50$), so the claim is
\emph{no measurable damage}, not improvement.

\subsection{Transfer}
\textbf{Across vocabularies.} Re-running the whole sweep against GPT-2's vocabulary
with the same settings: reduction percentages land within 0.1\,pp and the module
ranking is identical. The \emph{explanations} do worse --- one of our two stated reasons
for a tokenizer effect turned out to be true only of Qwen.

\textbf{Across conversations.} M7 mined from one half and tested on the other half:

\begin{center}
\begin{tabular}{@{}rrrr@{}}
\toprule
\textbf{Repetition in traffic} & \textbf{Gain} & \textbf{Extra cache hits} & \textbf{Extra gate refusals} \\
\midrule
0.0\%  & \textbf{+0.00\,pp}  & 0 & 0 \\
22.5\% & +5.43\,pp  & +4 & 0 \\
57.5\% & \textbf{+17.83\,pp} & +12 & 0 \\
\bottomrule
\end{tabular}
\end{center}

The exact zero at 0\% repetition is what makes the rest believable.

\subsection{Summary of metrics}
\begin{center}
\begin{tabularx}{\textwidth}{@{}X ll@{}}
\toprule
\textbf{Metric} & \textbf{Baseline} & \textbf{With Parsimony} \\
\midrule
Total token reduction & --- & \textbf{33.9\%} \\
Accuracy on 40 test questions & 92.5\% & \textbf{97.5\%}, no answer made worse \\
Wrong-answer rate on attack pairs & 26.7\% & \textbf{0.0\%} \\
Prefix reuse & 0.5\% & \textbf{98.5\%} \\
Time saved across the corpus & --- & \textbf{100.6\,s} (383\,ms per request) \\
Middleware overhead per request & --- & under 120\,ms \\
Mined-policy gain at 57\% repetition & --- & \textbf{+17.83\,pp} \\
\bottomrule
\end{tabularx}
\end{center}

\subsection{Threats to validity}
\begin{enumerate}
\item \textbf{One machine.} The 8.5\,ms per token figure is not claimed to hold elsewhere.
\item \textbf{Small test set.} 40 questions; after a grading fix only two defeat both models, and both are maths already handled without the model.
\item \textbf{Simple encoder.} Operating points move with the encoder; the ranking does not.
\item \textbf{Our corpus repeats only 1.9\%}, which is why \S6.5 reports a curve rather than one number. The repetition is synthetic; every question and answer in it is real.
\end{enumerate}

% ------------------------------------------------------------ conclusion
\section{Conclusion}

The efficiency literature has produced a rich toolbox and a thin account of what
happens when the tools are used together. Our contribution is not a new compression
algorithm but a \textbf{measuring instrument}: eight techniques in one harness, each
with its own switch, every decision written to a log that can be audited.

What that instrument found, again and again, is that \textbf{published settings are
specific to the setup they were tuned on}. A threshold called safe is unsafe here. A
placement rule backed by a well-cited paper is 80$\times$ more expensive here. A
calibration transfers as a ratio but not as an explanation. A routing assumption that
holds at 70B does not hold at 3B.

For this kind of system, the deliverable is \textbf{a calibration procedure, not a
number}.

% ------------------------------------------------------------ references
\begin{thebibliography}{40}
\footnotesize
\bibitem{r1} Li, Y., Dong, B., Lin, C., Guerin, F. \emph{Compressing Context to Enhance Inference Efficiency of LLMs.} EMNLP 2023. arXiv:2310.06201
\bibitem{r2} Jiang, H. et al. \emph{LLMLingua: Compressing Prompts for Accelerated Inference of LLMs.} EMNLP 2023. arXiv:2310.05736
\bibitem{r3} Jiang, H. et al. \emph{LongLLMLingua.} arXiv:2310.06839
\bibitem{r4} \emph{LLMLingua-2: Data Distillation for Task-Agnostic Prompt Compression.} arXiv:2403.12968
\bibitem{r5} \emph{Efficient Prompt Compression with Evaluator Heads.} arXiv:2501.12959
\bibitem{r6} \emph{SCOPE: A Generative Approach for LLM Prompt Compression.} arXiv:2508.15813
\bibitem{r7} \emph{Long Context In-Context Compression by Getting to the Gist of Gisting.} arXiv:2504.08934
\bibitem{r8} \emph{Compressing Lengthy Context With UltraGist.} arXiv:2405.16635
\bibitem{r9} \emph{ATACompressor: Adaptive Task-Aware Compression.} arXiv:2602.03226
\bibitem{r10} \emph{SARA: Selective and Adaptive Retrieval-augmented Generation.} arXiv:2507.05633
\bibitem{r11} \emph{PCToolkit: A Unified Prompt Compression Toolkit.} arXiv:2403.17411
\bibitem{r12} \emph{An Empirical Study on Prompt Compression for LLMs.} ICLR 2025.
\bibitem{r13} \emph{Understanding and Improving Information Preservation in Prompt Compression.} arXiv:2503.19114
\bibitem{r14} \emph{When Summaries Distort Decisions.} arXiv:2606.29251
\bibitem{r15} \emph{The Compression Paradox in LLM Inference.} arXiv:2603.23528
\bibitem{r16} Bang, F. \emph{GPTCache: An Open-Source Semantic Cache for LLM Applications.} NLP-OSS @ EMNLP 2023.
\bibitem{r17} \emph{GPT Semantic Cache.} arXiv:2411.05276
\bibitem{r18} \emph{MeanCache: User-Centric Semantic Caching for LLM Web Services.} arXiv:2403.02694
\bibitem{r19} \emph{ContextCache: Context-Aware Semantic Cache for Multi-Turn Queries.} arXiv:2506.22791
\bibitem{r20} \emph{A Generative Caching System for Large Language Models.} arXiv:2503.17603
\bibitem{r21} \emph{From Similarity to Vulnerability: Key Collision Attack on LLM Semantic Caching.} arXiv:2601.23088
\bibitem{r22} \emph{Enhancing adversarial resilience in semantic caching.} Scientific Reports, 2026.
\bibitem{r23} Kwon, W. et al. \emph{Efficient Memory Management for LLM Serving with PagedAttention.} SOSP 2023.
\bibitem{r24} vLLM. \emph{Automatic Prefix Caching.} docs.vllm.ai
\bibitem{r25} \emph{ChunkAttention: Prefix-Aware KV Cache and Two-Phase Partition.} arXiv:2402.15220
\bibitem{r26} \emph{Sparse Prefix Caching for Hybrid and Recurrent LLM Serving.} arXiv:2605.05219
\bibitem{r27} \emph{Multi-Segment Attention: Efficient KV-Cache Management.} arXiv:2606.02964
\bibitem{r28} Liu, N. F. et al. \emph{Lost in the Middle: How Language Models Use Long Contexts.} TACL 12:157--173, 2024.
\bibitem{r29} \emph{Adaptive Focus Memory for Language Models.} arXiv:2511.12712
\bibitem{r30} \emph{A Survey on Multi-Turn Interaction Capabilities of LLMs.} arXiv:2501.09959
\bibitem{r31} \emph{From Human Memory to AI Memory.} arXiv:2504.15965
\bibitem{r32} Chen, L., Zaharia, M., Zou, J. \emph{FrugalGPT.} arXiv:2305.05176
\bibitem{r33} Ong, I. et al. \emph{RouteLLM: Learning to Route LLMs with Preference Data.} arXiv:2406.18665
\bibitem{r34} \emph{When to Reason: Semantic Router for vLLM.} arXiv:2510.08731
\bibitem{r35} \emph{UCCI: Calibrated Uncertainty for Cost-Optimal LLM Cascade Routing.} arXiv:2605.18796
\bibitem{r36} \emph{Confident or Seek Stronger: Uncertainty-Based On-device LLM Routing.} arXiv:2502.04428
\bibitem{r37} \emph{Cluster, Route, Escalate: Cascaded Framework for Cost-Aware LLM Serving.} arXiv:2606.27457
\bibitem{r38} \emph{Precise length control for large language models.} Natural Language Processing Journal, Elsevier.
\bibitem{r39} \emph{BudgetThinker: Budget-aware LLM Reasoning with Control Tokens.} arXiv:2508.17196
\bibitem{r40} \emph{An Empirical Study of LLM Reasoning Ability Under Strict Output Length Constraint.} arXiv:2504.14350
\end{thebibliography}

\end{document}
